\rok{Фебруар 2026.}{B}

Други практични тест из Алгоритама и структура података

Први практични тест одржан 10.02.2026.

\zadatak{1}{Конверзија матрице у листу суседства}
Написати алгоритам који ће вршити директну конверзију матрице суседства у листу суседства. Алгоритам као улазни параметар треба да прихвати матрицу, а треба да резултује креираним динамичким низом листи.

\textbf{Додатне напомене за решавање задатка:}
\begin{itemize}
	\item У template-у је закоментарисана метода \texttt{ConvertAdjacencyMatrixToAdjacencyList} која треба да садржи имплементацију конверзије матрице у листу суседства.
	\item У Main методи обезбедити начин за тестирање и проверу нове функционалности.
\end{itemize}

\zadatak{2}{Уклањање свих карактера који се појављују у другом string-у}
Написати методу која прима два string-а. Из првог string-а треба уклонити све карактере који се налазе у другом string-у користећи хеш табелу ради ефикасности.

\textbf{Додатни захтеви за решавање задатка:}
\begin{itemize}
	\item Користити Hash табелу из шаблона за чување броја појављивања сваког карактера.
	\item Имплементацију писати у методи:
	\begin{Verbatim}[fontsize=\small]
public static void removePattern(string str, string pattern)
	\end{Verbatim}
	\item Временска комплексност треба да буде $O(n + m)$ где су „n“ и „m“ дужине прослеђених подstring-ова.
	\item Игнорисати разлику између малих и великих слова.
\end{itemize}

\textbf{Пример:}
\begin{itemize}
	\item \textbf{Улаз 1:} \texttt{str = "Algoritmi i strukture", pattern = "aeiou"}
	\item \textbf{Излаз 1:} \texttt{"lgrtm  strktr"} (уклоњени сви самогласници)
\end{itemize}

\zadatak{3}{Провера постојања путање са одређеном сумом}
Написати методу која проверава да ли у стаблу постоји путања од корена до било ког листа таква да сума вредности чворова на тој путањи износи тачно \texttt{targetSum}.

\textbf{Додатни захтеви за решавање задатка:}
\begin{itemize}
	\item Јавна метода:
	\begin{Verbatim}[fontsize=\small]
public boolean hasPathSum(int targetSum)
	\end{Verbatim}
	\item Приватна метода:
	\begin{Verbatim}[fontsize=\small]
private boolean hasPathSum(Node current, int currentSum, int targetSum)
	\end{Verbatim}
	\item Подразумева се да стабло има барем два чвора.
	\item Алгоритам ОБАВЕЗНО писати у оквиру класе \texttt{BinarySearchTree}.
\end{itemize}

\textbf{Примери:}

\begin{center}
\begin{minipage}{\textwidth}
\centering
\begin{forest}
for tree={
	circle,
	draw,
	thick,
	fill=gray!20,
	minimum size=8mm,
	inner sep=1pt,
	s sep=8mm,
	l sep=12mm,
	edge={-, thick},
	font=\small
}
[8
	[4
		[2
			[, phantom]
			[3]
		]
		[5]
	]
	[9
		[, phantom]
		[11]
	]
]
\end{forest}

\vspace{0.3cm}
\textbf{Слика 1.} Пример BST-a.
\end{minipage}
\end{center}

\begin{itemize}
	\item \textbf{Улаз 1:} \texttt{targetSum = 17}
	\item \textbf{Излаз 1:} \texttt{true} ($8 - 4 - 5$)
\end{itemize}

\begin{center}
\begin{minipage}{\textwidth}
\centering
\begin{forest}
for tree={
	circle,
	draw,
	thick,
	fill=gray!20,
	minimum size=8mm,
	inner sep=1pt,
	s sep=8mm,
	l sep=12mm,
	edge={-, thick},
	font=\small
}
[10
	[4
		[2
			[3]
			[, phantom]
		]
		[6
			[5]
			[9]
		]
	]
	[20
		[12
			[, phantom]
			[14]
		]
		[29
			[27]
			[, phantom]
		]
	]
]
\end{forest}

\vspace{0.3cm}
\textbf{Слика 2.} Пример BST-a.
\end{minipage}
\end{center}

\begin{itemize}
	\item \textbf{Улаз 2:} \texttt{targetSum = 99}
	\item \textbf{Излаз 2:} \texttt{false}
\end{itemize}

\sablon{Шаблон другог теста фебруар 2026. група Б}{2025/Februar/GrupaB/AISP2025_Test2_GrupaB.linq}
